% SOURCE_STATEMENT: "For general sufficiently small perturbations, possibly without symmetry, of generalized Kasner solutions of the Einstein equations, what behavior occurs near the singularity when the sub-criticality condition $\\max_{I<J,B}(\\widetilde q_I+\\widetilde q_J-\\widetilde q_B)<1$ fails?"
% Preparation and final self-audit: CONTRIBUTING.md prompts C and D; source/audit artifacts retained in the local physics_statements_2026_10_11 batch.
% !TeX program = lualatex
\documentclass[11pt]{article}
\usepackage[margin=25mm]{geometry}
\usepackage{fontspec}
\setmainfont{Latin Modern Roman}
\usepackage{amsmath,amssymb,mathtools}
\usepackage{unicode-math}
\setmathfont{Latin Modern Math}
\usepackage{xurl,hyperref}
\hypersetup{colorlinks=true,urlcolor=blue}
\setcounter{secnumdepth}{0}
\setlength{\parindent}{0pt}
\setlength{\parskip}{5pt}
\setlength{\emergencystretch}{3em}
\newcommand{\sref}[2]{\hyperlink{source:#1:#2}{[S#2]}}
\newcommand{\cataloguescope}[1]{\paragraph{Recorded target.}#1}
\newcommand{\R}{\mathbb R}
\newcommand{\C}{\mathbb C}
\newcommand{\N}{\mathbb N}
\newcommand{\Z}{\mathbb Z}
\newcommand{\Q}{\mathbb Q}
\newcommand{\F}{\mathbb F}
\newcommand{\eps}{\varepsilon}
\begin{document}
\section{Kasner Singularities Beyond the Subcritical Regime}\label{statement}

\subsection{Definitions and mathematical statement}

\paragraph{Shared notation.}$\mathbb N=\{1,2,\ldots\}$; the standard number fields have their usual meanings. All additional conventions are stated locally.


A generalized Kasner background has spatial scale factors $t^{\widetilde q_I}$ as proper time $t$ approaches zero. The exponents are those of the vacuum or Einstein--scalar-field background in the report. Asymptotically velocity-term-dominated behavior means that spatial derivative terms become negligible relative to time evolution. The subcritical condition is the displayed strict inequality in the target.

\paragraph{Catalogue wording retained for traceability.}
For general sufficiently small perturbations, possibly without symmetry, of generalized Kasner solutions of the Einstein equations, what behavior occurs near the singularity when the sub-criticality condition $\max_{I<J,B}(\widetilde q_I+\widetilde q_J-\widetilde q_B)<1$ fails?

\paragraph{Scope and source qualification.}The failure of the strict inequality includes the borderline equality case. Oscillatory behavior is a heuristic possibility, rather than an added asserted conclusion.
\subsection{Short English statement}
Determine the singular behavior of small unrestricted perturbations of Kasner solutions outside the proven subcritical regime.
\subsection{Sources}
\begin{itemize}
\item[\hypertarget{source:statement:1}{S1}] ULAM.AI and the UnsolvedMath Contributors, UnsolvedMath, catalogue entry. Catalogue curation under CC BY 4.0; definitions and English exposition added locally.
\url{https://huggingface.co/datasets/ulamai/UnsolvedMath}.
\item[\hypertarget{source:statement:2}{S2}] Mathematical Aspects of General Relativity (2021), Oberwolfach Reports, DOI 10.4171/owr/2021/40, pp. 2157--2267. Original report and mathematical context.
\url{https://doi.org/10.4171/owr/2021/40}.
\end{itemize}
\label{end:statement}
\end{document}
